\section{Concepts D: Data-Level Parallelism (Vector, SIMD, GPU/SIMT)}\label{grp:dlp}
\textit{Textbook (H\&P Ch.~4); not in the slides but named by the professor.} \textbf{Why DLP is efficient:} one instruction controls $N$ operations, so fetch/decode/dependence checks are paid once per $N$; elements are independent by definition; the access pattern is known. \textbf{Spot it:} iteration $i$ touches only element $i$ (no order between elements). A loop-carried dependence blocks it, except \emph{reductions} (tree, $\log_2N$ steps) and scans. Write: ``DLP present because\ldots; the only ordering is the final reduction.''
{\small\renewcommand{\arraystretch}{1.3}\rowcolors{1}{white}{white}
\begin{longtable}{L{2.3cm}L{4.7cm}L{2.9cm}L{3.0cm}L{2.8cm}}
\toprule
\rowcolor{hdrc}\textbf{Concept} & \textbf{What it fixes $\cdot$ how it works} & \textbf{Best scenario (use when)} & \textbf{Why it beats the alternatives} & \textbf{Weakness $\cdot$ avoid when}\\
\midrule
\endfirsthead
\toprule
\rowcolor{hdrc}\textbf{Concept} & \textbf{What it fixes $\cdot$ how it works} & \textbf{Best scenario (use when)} & \textbf{Why it beats the alternatives} & \textbf{Weakness $\cdot$ avoid when}\\
\midrule
\endhead
\bottomrule
\endfoot
\textbf{\textcolor{cardc}{D1}}\label{card:D1}\newline\textbf{Vector architecture}\newline{\footnotesize\textit{Static (code) + hardware}} & \textbf{Fixes:} Long regular vectors needing high predictable throughput. \newline \textbf{How:} Vector registers (VMIPS 8$\times$64$\times$64-bit); VL register for $n<$MVL; \emph{strip-mining} for $n>$MVL; \emph{lanes} (element $i\rightarrow$ lane $i \bmod L$); \emph{chaining} (vector forwarding); \emph{mask registers} for \texttt{if}; strided loads; gather/scatter for sparse. & Scientific/supercomputer kernels with long vectors. & SIMD extensions: variable length, compact code, hardware strides/gather. Multicore: far less control per operation. & Needs many memory banks; poor on short vectors; gather/scatter slow. \\*
\rowcolor{black!4}\multicolumn{5}{p{\dimexpr\linewidth-2\tabcolsep\relax}}{\textbf{Numbers:} Time $\approx$ convoys $\times\ n/L$ + start-up. Banks $\ge$ bank busy time; stride $s$ conflicts if banks$/\gcd(s,\text{banks})<$ busy time. \enspace$\cdot$\enspace \textbf{\textcolor{pickc}{Write:}} \textit{Each vector instruction does [VL] operations across [L] lanes with chaining, so [n] elements take about [convoys]$\times n/L$ cycles with one instruction fetch per vector.}}\\
\midrule
\textbf{\textcolor{cardc}{D2}}\label{card:D2}\newline\textbf{SIMD extensions on the CPU (SSE, AVX, NEON)}\newline{\footnotesize\textit{Static}} & \textbf{Fixes:} DLP in short vectors without leaving the CPU. \newline \textbf{How:} 128/256/512-bit registers of packed elements (8 floats in 256 bits); compiler or intrinsics vectorise loops. & Media, crypto, small/medium kernels, data already in CPU caches, no transfer wanted. & GPU for small problems: no PCIe transfer or launch overhead. & Width fixed in the ISA; alignment; weak gather/masks (better in AVX-512); speedup $\le$ elements per register; compilers vectorise only simple loops. \\*
\rowcolor{black!4}\multicolumn{5}{p{\dimexpr\linewidth-2\tabcolsep\relax}}{\textbf{Needs:} Memory bandwidth to feed it \cp{M18}; unit-stride/SoA layout \cp{M13}. \enspace$\cdot$\enspace \textbf{\textcolor{pickc}{Write:}} \textit{The loop applies one operation to [n] contiguous floats; 256-bit SIMD processes 8 per instruction in the existing core, with no PCIe transfer.}}\\
\midrule
\textbf{\textcolor{cardc}{D3}}\label{card:D3}\newline\textbf{GPU / SIMT}\newline{\footnotesize\textit{Dynamic (warp scheduling)}} & \textbf{Fixes:} Massive DLP; throughput over latency. \newline \textbf{How:} Programmer writes one scalar thread; grid of thread blocks; each block on one SM; 32 threads form a \textbf{warp} executing one instruction across 32 lanes. Memory latency hidden by switching to another ready warp each cycle (fine-grained MT, all warps' registers resident), not by big caches/OoO. & $10^5$+ independent elements, regular arithmetic, data resident on the device. & Multicore CPU: thousands of lanes vs.\ tens; more ops per watt; huge memory bandwidth. & Small caches per thread; needs parallel slack; transfer and programming cost. \newline \textbf{Avoid:} Branchy per-element code, pointer structures, small problems, data that must cross PCIe for little work. \\*
\rowcolor{black!4}\multicolumn{5}{p{\dimexpr\linewidth-2\tabcolsep\relax}}{\textbf{Needs:} Occupancy \cp{D8}, coalescing \cp{D4}, low divergence \cp{D5}, transfer plan \cp{D7}. \enspace$\cdot$\enspace \textbf{\textcolor{pickc}{Write:}} \textit{[$10^6$] independent pixels give enough warps to hide [400]-cycle memory latency by warp switching, so a GPU's thousands of lanes stay busy where an 8-core CPU has only tens.}}\\
\midrule
\textbf{\textcolor{cardc}{D4}}\label{card:D4}\newline\textbf{Coalesced access / data layout (SoA)}\newline{\footnotesize\textit{Static}} & \textbf{Fixes:} Uncoalesced memory: up to 32 transactions per warp instead of one. \newline \textbf{How:} Thread $i$ accesses element $i$ (consecutive addresses); struct-of-arrays; stage awkward patterns (transpose) through shared memory. & --- & More bandwidth/bigger GPU: waste scales with them. & --- \\*
\rowcolor{black!4}\multicolumn{5}{p{\dimexpr\linewidth-2\tabcolsep\relax}}{\textbf{Numbers:} Row-per-thread on a 3840-wide RGBA image: neighbours in a warp are 15\,KB apart $\Rightarrow$ 32 transactions. \enspace$\cdot$\enspace \textbf{\textcolor{pickc}{Write:}} \textit{Mapping thread $i$ to element $i$ makes a warp read 32 consecutive words in one transaction instead of 32 separate ones.}}\\
\midrule
\textbf{\textcolor{cardc}{D5}}\label{card:D5}\newline\textbf{Divergence handling (compaction, sorting, predication)}\newline{\footnotesize\textit{Dynamic / Static}} & \textbf{Fixes:} Warps executing both sides of a data-dependent \texttt{if}. \newline \textbf{How:} Hardware masks lanes and runs each side in turn, then reconverges. Fix: \emph{compaction} (pass 1 lists elements needing the slow path; pass 2 runs them in full warps: dynamic); \emph{sort/bin} data so a warp branches alike; \emph{predicate} short branches (static). & --- & Predication alone (pays both sides); bigger GPU (same waste). & Extra pass and list memory; atomics on the list counter (aggregate per warp). \\*
\rowcolor{black!4}\multicolumn{5}{p{\dimexpr\linewidth-2\tabcolsep\relax}}{\textbf{Numbers:} Efficiency $=$ active lanes/32. 30\% random slow pixels (50 vs.\ 5 instr.): $P(\text{no slow in warp})=0.7^{32}\approx0$ $\Rightarrow$ 55 executed vs.\ 18.5 useful $\Rightarrow$ 34\%. \enspace$\cdot$\enspace \textbf{\textcolor{pickc}{Write:}} \textit{Because [30\%] of elements, randomly placed, take the slow path, nearly every warp diverges (efficiency [34\%]); compacting them into a separate list gives full, uniform warps.}}\\
\midrule
\textbf{\textcolor{cardc}{D6}}\label{card:D6}\newline\textbf{Shared memory (scratchpad) tiling}\newline{\footnotesize\textit{Static}} & \textbf{Fixes:} Neighbourhood reuse re-reading global memory. \newline \textbf{How:} A block loads its tile (+ halo) once into on-chip shared memory, then all threads reuse it; barrier inside the block after the load. & Stencils, filters, matrix multiply; predictable timing (real-time). & Hardware cache: software decides what stays $\Rightarrow$ no conflict/eviction surprises; raises arithmetic intensity. & Programmer effort; limits occupancy if too large. \\*
\rowcolor{black!4}\multicolumn{5}{p{\dimexpr\linewidth-2\tabcolsep\relax}}{\textbf{\textcolor{pickc}{Write:}} \textit{Each output reads a [$3\times3$] neighbourhood; loading the tile once into shared memory serves 9 reads from on-chip storage.}}\\
\midrule
\textbf{\textcolor{cardc}{D7}}\label{card:D7}\newline\textbf{Host--device transfer: residency, batching, overlap}\newline{\footnotesize\textit{Static}} & \textbf{Fixes:} PCIe copies dominating kernel time. \newline \textbf{How:} Keep data on the GPU across kernels; batch small jobs; double-buffer (copy frame $k{+}1$ while computing $k$). & --- & --- & \textbf{Avoid:} If compute per byte is tiny: keep it on CPU SIMD \cp{D2}. \\*
\rowcolor{black!4}\multicolumn{5}{p{\dimexpr\linewidth-2\tabcolsep\relax}}{\textbf{Numbers:} 33\,MB frame over 16\,GB/s $\approx$2\,ms each way. \enspace$\cdot$\enspace \textbf{\textcolor{pickc}{Write:}} \textit{Copies take [4] of every [10]\,ms; keeping data resident and double-buffering overlaps them with compute.}}\\
\midrule
\textbf{\textcolor{cardc}{D8}}\label{card:D8}\newline\textbf{Occupancy tuning}\newline{\footnotesize\textit{Static}} & \textbf{Fixes:} Too few resident warps to hide memory latency. \newline \textbf{How:} Occupancy $=$ resident warps / max per SM; raise it with more threads (one per element), fewer registers per thread, moderate shared memory, block size 128--256\bs. & --- & --- & More registers or shared memory per thread lowers the number of resident warps. \\*
\rowcolor{black!4}\multicolumn{5}{p{\dimexpr\linewidth-2\tabcolsep\relax}}{\textbf{Numbers:} 2160 threads on 40 SMs $=68$ warps ($<$2 per SM): latency exposed. \enspace$\cdot$\enspace \textbf{\textcolor{pickc}{Write:}} \textit{Only [68] warps for [40] SMs cannot hide latency; one thread per element gives [thousands] of warps.}}\\
\end{longtable}}

\small
\begin{tabularx}{\textwidth}{L{2.4cm}L{3.2cm}L{3.2cm}L{3.6cm}X}
\toprule
\tblhead{\textbf{} & \textbf{Vector} & \textbf{SIMD extension} & \textbf{GPU (SIMT)} & \textbf{Multicore CPU (MIMD)}}\\
\midrule
Control & 1 instr, VL elements & 1 instr, fixed width & 1 instr per 32-thread warp & 1 PC per core \\
Hides latency by & Deep pipelines, banks & Caches, OoO & Switching warps & Caches, OoO, SMT \\
Branches & Masks & Masks/blend & Masks + reconvergence (divergence) & Free \\
Irregular memory & Gather/scatter & Weak & Uncoalesced = slow & Caches tolerate it \\
Best for & Long regular vectors & Short vectors, no transfer & Huge regular throughput & Irregular, branchy, task-parallel \\
Weak at & Short vectors & ISA width churn & Transfer, divergence, small jobs & Ops per watt on regular data \\
\bottomrule
\end{tabularx}
\normalsize
\begin{trap}[DLP traps]
A GPU often wins by \emph{memory bandwidth}, not FLOPs: SAXPY AI $=2/12$ FLOP/B is memory-bound everywhere (roofline \pg{sec:roofline}). More lanes do not fix divergence or uncoalesced access. Loop-carried dependences cannot be vectorised (except reductions/scans). GPUs need parallel slack: too little work runs badly.
\end{trap}
